\section{Score Matching en tiempo continuo}

\subsection{Objetivo de entrenamiento}

En el marco de tiempo continuo, el objetivo de entrenamiento del score matching se puede expresar de manera unificada como:

$$
\mathcal{L} = \mathbb{E}_{t \sim \mathcal{U}(0,T)}\;\mathbb{E}_{\mathbf{x}_0 \sim p_{\text{data}}}\;\mathbb{E}_{\mathbf{x}_t \sim q(\mathbf{x}_t|\mathbf{x}_0)}\left[\lambda(t)\left\|\mathbf{s}_\theta(\mathbf{x}_t, t) - \nabla_{\mathbf{x}_t}\log q(\mathbf{x}_t|\mathbf{x}_0)\right\|^2\right]
$$

Donde:
\begin{itemize}
    \item $t$ se muestra de forma uniforme en $[0, T]$ (tiempo continuo);
    \item $\lambda(t)$ es una función de ponderación dependiente del tiempo;
    \item El score condicional $\nabla_{\mathbf{x}_t}\log q(\mathbf{x}_t|\mathbf{x}_0) = -\frac{\mathbf{x}_t - \alpha_t \mathbf{x}_0}{\sigma_t^2}$ tiene una forma analítica conocida.
\end{itemize}

\begin{importantbox}{Entrenamiento discreto vs. continuo}
En el DDPM discreto, el paso de tiempo $t$ se muestra uniformemente desde $\{1, 2, \ldots, T\}$. En el marco continuo, $t$ se muestra desde un intervalo continuo $[0, T]$. En la práctica del entrenamiento, ambas diferencias son prácticamente insignificantes: basta con reemplazar los índices de tiempo discretos por valores continuos de tiempo. La red $\mathbf{s}_\theta(\mathbf{x}, t)$ acepta $t$ continuo como entrada (normalmente procesado mediante codificación posicional sinusoidal).
\end{importantbox}

\subsection{Selección de la función de ponderación $\lambda(t)$}

Diferentes elecciones de $\lambda(t)$ corresponden a distintas estrategias de entrenamiento:

\begin{itemize}
    \item $\lambda(t) = g(t)^2$: corresponde al ponderado de verosimilitud en las EDE, relacionado con la optimización del límite variacional (VLB);
    \item $\lambda(t) = \sigma_t^2$: es la elección más común; equivale al objetivo de predicción de ruido $\|\boldsymbol{\epsilon}_\theta - \boldsymbol{\epsilon}\|^2$;
    \item $\lambda(t) = 1$: coincide directamente el score, pero puede ser inestable en momentos de bajo ruido (pequeños $t$).
\end{itemize}

\begin{knowledgebox}{De Score Matching a predicción de ruido}
Utilizando $\nabla_{\mathbf{x}_t}\log q(\mathbf{x}_t|\mathbf{x}_0) = -\boldsymbol{\epsilon}/\sigma_t$ y $\mathbf{s}_\theta = -\boldsymbol{\epsilon}_\theta / \sigma_t$, la pérdida de score matching se puede reescribir como:
$$
\left\|\mathbf{s}_\theta - \nabla_{\mathbf{x}_t}\log q\right\|^2 = \frac{1}{\sigma_t^2}\left\|\boldsymbol{\epsilon}_\theta - \boldsymbol{\epsilon}\right\|^2
$$
Por lo tanto, elegir $\lambda(t) = \sigma_t^2$ elimina exactamente el factor $1/\sigma_t^2$, obteniendo una pérdida de mínimos cuadrados medios (MSE) de predicción de ruido simplificada: este es precisamente el objetivo simplificado utilizado en el artículo original de DDPM.
\end{knowledgebox}

\subsection{Resumen de la sección}

El objetivo de entrenamiento del score matching en el marco de tiempo continuo es sustancialmente idéntico al del DDPM discreto; la única diferencia radica en que la variable de tiempo pasa de ser discreta a continua. La elección de la función de ponderación $\lambda(t)$ influye en la dinámica del entrenamiento, pero el objetivo más frecuente de predicción de ruido es completamente consistente desde ambas perspectivas.
